\section{Brownian diffusion and the escape observable}
\label{sec:diffusion}

\subsection{From Fick's law to a lattice random walk}

In a homogeneous fluid without advection or external forces, conservation of
mass together with Fick's first law, $\bm J=-D\bm\nabla c$, gives Fick's second
law
\begin{equation}
 \frac{\partial c}{\partial t}=D\nabla^2c.
 \label{eq:fick2}
\end{equation}
The forward-Euler, centered-space discretization of Eq.~\ref{eq:fick2} at
lattice site $\bm i$ is
\begin{equation}
 c_{\bm i}^{n+1}
 =c_{\bm i}^{n}
 +\lambda_D\sum_{\bm e\in\mathcal N_6}
 \left(c_{\bm i+\bm e}^{n}-c_{\bm i}^{n}\right),
 \label{eq:ftcs}
\end{equation}
where $\mathcal N_6=\{\pm\hat{\bm x},\pm\hat{\bm y},\pm\hat{\bm z}\}$.
Equation~\ref{eq:ftcs} is also the master equation of the random-walk kernel
defined by Eqs.~\ref{eq:lambda-D} and \ref{eq:stay-probability}.  Thus, the
single-particle Monte Carlo propagator and the finite-difference diffusion
equation have the same ensemble evolution.

For one update, symmetry gives $\E[\Delta\bm r]=\bm 0$ and
\begin{equation}
 \E\!\left[\lvert\Delta\bm r\rvert^2\right]
 =6\lambda_Da^2=6D\dt.
 \label{eq:one-step-msd}
\end{equation}
Independent increments therefore give
\begin{equation}
 \E\!\left[\lvert\bm r(t)-\bm r(0)\rvert^2\right]=6Dt,
 \qquad
 \E[(x(t)-x(0))^2]=2Dt,
 \label{eq:einstein-msd}
\end{equation}
which is the Einstein relation for three-dimensional Brownian diffusion
\citep{einstein1905}.  The lattice walk consequently represents overdamped
Brownian transport on times and lengths resolved by $\dt$ and $a$.

\subsection{Boundary conditions}

The manuscript simulations use periodic lateral boundaries,
\begin{equation}
 (x,y,z)\equiv(x+L_x,y,z)\equiv(x,y+L_y,z),
 \label{eq:periodic-boundaries}
\end{equation}
implemented by modular wrapping of the $x$ and $y$ indices.  This suppresses
artificial loss from the sides and represents a periodically tiled sensor
patch.  A move proposed into the voxelized solid is rejected, which is the
random-walk realization of a reflecting, no-flux boundary on nonreactive
surface encounters.  Receptors on exposed faces add stochastic reaction to
this otherwise reflecting interface.

For the escape assay, a trajectory terminates when the free ligand reaches an
absorbing escape surface.  The principal manuscript calculations use a
horizontal plane $z=z_{\mathrm{esc}}$; crossing an explicitly open domain
boundary also counts as escape.  If
\begin{equation}
 T_{\mathrm{esc}}=\inf\{t>0:z(t)\ge z_{\mathrm{esc}}\},
 \label{eq:first-passage}
\end{equation}
then the retention survival function is the first-passage quantity
\begin{equation}
 S(t)=\Prb(T_{\mathrm{esc}}>t).
 \label{eq:survival}
\end{equation}
It measures persistence in the explicit near-surface region, not the survival
of a chemical bond.  A trajectory can alternate many times between bound and
unbound states before the event in Eq.~\ref{eq:first-passage}.

Trajectories that have not escaped by the finite horizon $t_{\max}$ are
right-censored.  With $d_j$ escape events among $n_j$ trajectories at risk at
ordered event time $t_j$, the plotted Kaplan--Meier estimator is
\begin{equation}
 \widehat S(t)=\prod_{t_j\le t}\left(1-\frac{d_j}{n_j}\right).
 \label{eq:kaplan-meier}
\end{equation}
The assay additionally records the observation time, number of binding
intervals, self- and cross-rebinding counts, and cumulative bound time
\begin{equation}
 T_B=\sum_{q=0}^{N_{\mathrm{reb}}}B_q,
 \label{eq:cumulative-bound-time}
\end{equation}
where $B_0$ is the initial bound interval and subsequent $B_q$ are rebound
intervals.  Free-excursion durations contribute to $T_{\mathrm{esc}}$ but not
to $T_B$.  This distinction is essential because receptor occupancy is a
time-integrated bound-state observable.
